\documentclass[12pt,a4paper]{article}

% ------------------------
% Paquetes
% ------------------------
\usepackage[spanish]{babel}
\usepackage[utf8]{inputenc}
\usepackage[T1]{fontenc}
\usepackage{lmodern}
\usepackage{geometry}
\usepackage{enumitem}
\usepackage{booktabs}
\usepackage{hyperref}
\geometry{margin=2.5cm}

\begin{document}

% ------------------------
\section*{Learning Objectives}
\subsection*{General Objective}
The student will apply the analysis, design, development, and implementation of transactional and analytical information systems with structured data through readings, essays, in-class exercises, and practical work, so that by the end of the course they will be capable of creating OLTP and OLAP information systems.

% ------------------------
\section*{Course Content}
\begin{enumerate}
    \item \textbf{Introduction to Information Systems}
    \begin{enumerate}[label*=\arabic*.]
        \item Concept of Information System.
        \item Components:
        \begin{enumerate}[label*=\arabic*.]
            \item Users and processes.
            \item Hardware and networks, client–server architecture.
            \item Software.
            \item Operating Systems.
            \item Database Management System (DBMS).
            \item Security, confidentiality, backup, and recovery.
            \item Concurrency.
            \item Database.
            \item Integrity and consistency.
        \end{enumerate}
        \item Main Types of Information Systems:
        \begin{enumerate}[label*=\arabic*.]
            \item Characteristics of Transactional Systems.
            \item Characteristics of Analytical Systems.
        \end{enumerate}
        \item Transactional Information Systems (OLTP):
        \begin{enumerate}[label*=\arabic*.]
            \item Concept of a transaction.
            \item ACID properties of a transaction.
        \end{enumerate}
    \end{enumerate}
    \item \textbf{Analysis, Design, and Implementation of Transactional Information Systems (OLTP)}
    \begin{enumerate}[label*=\arabic*.]
        \item Requirements Analysis:
        \begin{enumerate}[label*=\arabic*.]
            \item Conceptual design.
            \item Entity–Relationship Model.
            \item Practical exercises.
        \end{enumerate}
        \item Logical Design:
        \begin{enumerate}[label*=\arabic*.]
            \item Relational Model.
            \item Normalization.
            \item Relational algebra.
        \end{enumerate}
        \item Physical Design.
        \item SQL:
        \begin{enumerate}[label*=\arabic*.]
            \item Data Definition Language (DDL).
            \item Data Control Language (DCL).
            \item Data Query Language (DQL).
            \item SQL query exercises.
            \item Codd's Twelve Rules.
            \item The relational and in-memory model as the best option for transactional systems.
        \end{enumerate}
    \end{enumerate}
    \item \textbf{Analysis, Design, and Implementation of Analytical Information Systems (OLAP)}
    \begin{enumerate}[label*=\arabic*.]
        \item Challenges and issues in Business Intelligence.
        \item Analysis and design of the data warehouse.
        \item Models: multidimensional, snowflake, star.
        \item Creating a data warehouse with SQL.
    \end{enumerate}
    \item \textbf{ETL Process}
    \begin{enumerate}[label*=\arabic*.]
        \item Extraction.
        \item Cleaning and standardization.
        \item Transformation.
        \item Loading.
        \item ETL tools.
    \end{enumerate}
    \item \textbf{Business Intelligence Components}
    \begin{enumerate}[label*=\arabic*.]
        \item Descriptive and predictive analytics.
        \item Complex SQL queries.
        \item Data mining.
        \item OLAP tools, dashboards, and scorecards.
    \end{enumerate}
\end{enumerate}

% ------------------------
\section*{Course Evaluation}
\begin{table}[h!]
\centering
\begin{tabular}{@{}ll@{}}
\toprule
\textbf{Activity} & \textbf{Percentage} \\
\midrule
Midterm Exams & 50\% \\
Participation & 10\% \\
Practical Work (Mandatory) & 20\% \\
Project (Mandatory) & 20\% \\
\midrule
\textbf{Total} & \textbf{100\%} \\
\bottomrule
\end{tabular}
\end{table}

\noindent \textit{Note}: In order to qualify for the first final exam, it is mandatory to submit at least 80\% of the practical work and projects.

% ------------------------
\section*{Basic Bibliography}
\begin{enumerate}
    \item Coronel, C., \& Morris, S. (2016). \textit{Database Systems}. Cengage Learning.
    \item Date, C. (2007). \textit{An Introduction to Database Systems}. Pearson Education.
    \item Inmon, W. (2011). \textit{Building the Data Warehouse}. Wiley.
    \item Kimball, R., \& Ross, M. (2013). \textit{The Data Warehouse Toolkit}. Wiley.
    \item Lans, R. (2012). \textit{Data Virtualization for Business Intelligence Systems}. Morgan Kaufmann.
    \item Ricardo, C., \& Urban, S. (2017). \textit{Databases Illuminated}. Jones \& Bartlett Learning.
\end{enumerate}

% ------------------------
\section*{Resources and Tools}
\begin{itemize}
    \item PowerDesigner Trial: \url{https://www.powerdesigner.biz/EN/powerdesigner/powerdesigner-trial.html}
    \item Bizagi BPMN Modeler: \url{https://www.bizagi.com/en/platform/modeler}
    \item AltairOne Registration: \url{https://admin.altairone.com/register}
    \item RapidMiner: \url{https://docs.rapidminer.com/latest/studio/installation}
    \item Qlik Sense Academic Program: \url{https://www.qlik.com/us/company/academic-program}
\end{itemize}

\section{Lecture I}
\subsection{From File Systems to Intelligent Databases: OLTP, OLAP, and Beyond}

In the early days of computing, information was stored in what we now call \emph{file systems}.  
A program had to know exactly where and how the data was written.  
Data and code were tightly coupled: change one, and you might break the other.  
There was no real notion of \textbf{integrity} or \textbf{consistency}, and no way to guarantee that a transaction --- a complete logical unit of work --- could succeed or fail as a whole.  

This world was known as \textbf{Standalone Applications (SA)}.  
The great leap came with \textbf{Database Systems (SBD)}, where data became independent of the code, and the database engine (the \emph{Database Management System}, or DBMS) took responsibility for enforcing structure, integrity, and transaction control.

\paragraph{Three Fundamental Concepts}
Before diving into history, let us clear up three often-confused terms:
\begin{enumerate}
    \item \textbf{Database (DB):} A structured collection of data, interconnected in such a way that it remains logically \emph{consistent} and \emph{integral}.
    \item \textbf{Database Management System (DBMS):} The software responsible for storing, securing, organizing, accessing, processing, and keeping the data integral and consistent. It handles queries, security, concurrency, and recovery.
    \item \textbf{Database System:} The complete information system where the data resides in a database, managed by a DBMS, running on an operating system, stored in hardware, and executed on CPUs, typically connected through a client–server architecture.
\end{enumerate}

\paragraph{On Integrity}
In everyday language, ``integrity'' means wholeness or incorruptibility.  
In databases, integrity is the property that the data \emph{represents reality faithfully} and \emph{conforms to defined rules}.  
A database that is integral enforces that each value respects its domain, its relationships, and its meaning.  
A field \texttt{birth\_date} containing \texttt{2050-12-01} violates temporal logic;  
a \texttt{name} field storing ``Ballenttina'' when the person’s real name is ``Valentina'' violates semantic integrity even though it may be syntactically valid.  
Thus, integrity goes beyond data type checks: it demands correctness in context.

\subsubsection{OLTP: Online Transaction Processing}
In the 1970s and 1980s, as organizations sought to automate daily operations, \textbf{OLTP systems} emerged.  
They excel at:
\begin{itemize}
    \item Handling \emph{many users} simultaneously.
    \item Managing \emph{a high volume of small transactions}.
    \item Performing mostly \emph{write} and \emph{update} operations on relatively small datasets.
\end{itemize}
Banking systems, airline reservation platforms, and retail point-of-sale systems are classic examples.  
With the advent of the Internet and digital communications, OLTP systems began handling \emph{semi-structured data} (e.g., XML, JSON).  
Response times moved into the realm of milliseconds and nanoseconds.  
Early OLTP was built with languages like \texttt{COBOL} and \texttt{C}, later enriched by \texttt{SQL}.

Yet, as OLTP matured, an uncomfortable truth emerged: even the most perfectly designed database is useless if the \emph{data quality} itself is poor.

\subsubsection{OLAP: Online Analytical Processing}
By the 1990s, businesses wanted more than just daily transaction logs: they wanted to \emph{understand} the past and guide the future.  
\textbf{OLAP systems} emerged to analyze large volumes of historical data, providing aggregated views, trend analyses, and strategic insights.  
Unlike OLTP:
\begin{itemize}
    \item They have \emph{few users}, often analysts or executives.
    \item They focus on \emph{read-only operations}.
    \item Response times are measured in \emph{minutes or hours}.
\end{itemize}
This era also saw the rise of \textbf{Business Intelligence} and the development of the \textbf{ETL} process (\emph{Extract, Transform, Load}), which moves data from OLTP systems into data warehouses.  
OLAP initially relied on \texttt{SQL} for relational querying, later expanding to multidimensional query languages like \texttt{MDX}.

\subsubsection{From OLAP to Data Science}
In the 21st century, data analysis expanded into what we now call \textbf{Data Science}, encompassing:
\begin{itemize}
    \item \textbf{Descriptive analytics} --- understanding the past and present.
    \item \textbf{Predictive analytics} --- anticipating likely future scenarios.
    \item \textbf{Prescriptive analytics} --- identifying actions today to alter tomorrow’s outcomes.
\end{itemize}
This shift absorbed structured, semi-structured, and unstructured data sources.  
Data volumes grew enormously, user bases became diverse, and operations focused mainly on reading, modeling, and extracting patterns.  
Still, one ancient problem remains unsolved: the quality of the data itself.

\subsection{Integrity, Consistency, and the OLTP Design Process}

When we speak of \emph{integrity} in databases, we often mean it in a loose, everyday sense:  
``the data is correct.'' But in database theory, integrity has a sharper structure, and we usually formalize it into three distinct forms.

\paragraph{Domain Integrity}  
Each attribute in a database has a \emph{domain} --- the set of values it can validly take.  
A column for a person’s date of birth should only contain valid dates;  
a column for an e-mail address should follow the correct syntactic pattern.  
Domain integrity enforces these rules at the type and format level.

\paragraph{Business Rule (Value) Integrity}  
Beyond technical types, there are the rules that arise from the meaning of the data.  
A business might declare that no employee can have more than one active contract,  
or that a customer’s credit limit cannot be negative.  
These are \emph{semantic} constraints, rooted in the logic of the enterprise itself.

\paragraph{Referential Integrity}  
This is the hallmark of \textbf{relational} databases.  
It ensures that relationships between tables remain valid:  
a foreign key in one table must match a primary key in another,  
or be null if the relationship allows it.  
If you delete a customer, you must also decide what happens to their orders ---  
do they vanish (\emph{cascade}), stay orphaned (\emph{set null}), or remain protected (\emph{restrict})?

\paragraph{On Redundancy and Consistency}  
Redundancy occurs when the same piece of information appears multiple times unnecessarily.  
Imagine two separate databases: one for human resources, listing employees and their contact details;  
another for customer management, also listing contact details for certain employees who are client liaisons.  
If an employee changes their phone number, both databases must be updated ---  
and if they are not, the result is a \emph{loss of consistency}.  
Consistency here means that, across all datasets, the representation of reality is aligned and coherent.

\subsubsection{OLTP Analysis and Design}

An \textbf{OLTP transaction} is a set of insert, delete, or update operations that together form an indivisible unit.  
Before and after the transaction, the database must remain integral and consistent.  
Inside the transaction, the operations are executed in sequence, invisible to the outside world until the transaction is complete.

\paragraph{Concurrency Control}  
A DBMS must allow multiple users to access and modify data at the same time without corrupting it.  
This is achieved through mechanisms such as scheduling (serializing transactions),  
and locks (ensuring that two conflicting operations cannot occur simultaneously).  
The DBMS enforces rules so that all users perceive the database as if their transactions were the only ones running.

\paragraph{The ACID Properties}  
Transactions are defined by four properties, captured by the mnemonic \textbf{ACID}:
\begin{enumerate}
    \item \textbf{Atomicity} --- All operations in a transaction are executed, or none are. There is no in-between.
    \item \textbf{Consistency} --- The transaction takes the database from one consistent state to another.  
    Even if the transaction fails, the system ensures that no inconsistent state is left behind.
    \item \textbf{Isolation} --- While a transaction is running, its intermediate results are invisible to other transactions.
    \item \textbf{Durability} --- Once a transaction is committed, its changes persist even in the face of system failures.  
    This is achieved through logs (before and after images of the data),  
    allowing recovery to the exact committed state.
\end{enumerate}

\subsubsection{Three Levels of OLTP Database Design}

Designing an OLTP system can be viewed as working through three abstraction levels:

\paragraph{1) Conceptual Modeling}  
We begin with \emph{requirements analysis}:  
collecting the needs of the system and identifying the entities (nouns) and relationships (verbs) involved.  
The output is the \textbf{Entity–Relationship Diagram} (ERD), which describes the logical structure of the data without concern for implementation.

\paragraph{2) Logical Modeling}  
The ERD is transformed into a \textbf{relational model}.  
This involves \emph{normalization}, a process that removes redundancy and ensures integrity by organizing data into well-structured tables.  
The output is a fully normalized relational schema.

\paragraph{3) Physical Modeling}  
Finally, the logical model is adapted to the specific DBMS and hardware environment.  
Here we decide on indexes, partitioning strategies, and storage structures,  
always aiming to maintain integrity, consistency, and performance.  
The output is the actual database and its elements, ready to handle live transactions.


What is the relational and the logical model? how does the peter chen ERM differs from ER Merise
El trabajo publicado por Codd en ACM (1970) presentaba un nuevo modelo de datos que perseguia una serie de objetivos, que se resumen en las siguientes lineas.
I.Fisica: El modo en el que se almacenan los datos no influye en sumanipulacion logica y por tanto, los usuarios que acceden a esos datos no tienen que modificar sus programas por cambios en el almacenamiento fisico.
I.Logica: El anadir, eliminar o modificar objetivos de la base da datos (es decir cambios en el esquema) no repercute en los programas y/o usuarios que estan consultando a subconjuntos parciales de los mismos (vistas)

Details of te relational model: attribute, tuple, relation, cardinality, degree, p key, f key, data type, esquema, instance, intension.

La tupla representa entidad
La relacion representa al conjunto de entidades

La intension es relacionada al esquema
La extension se relaciona con los datos



\end{document}

